\documentclass[../../main.tex]{subfiles}
\begin{document}
\section{子流形的切空间}\label{section:ch05-子流形的切空间}

设 $S\subseteq M$ 是浸入子流形,并令 $\iota:S\hookrightarrow M$ 为包含映射。由于 $\iota$ 是浸入, $\mathrm{d}\iota_p$ 是单射,故我们可把 $T_pS$ 视为 $T_pM$ 的子空间,只需把 $T_pS$ 等同于 $\mathrm{d}\iota_p(T_pS)\subseteq T_pM$。作此等同后, $T_pS$ 称为 $S$ 在 $p$ 处的\textbf{切空间}。由于 $\iota$ 是映射 $S\to M$ 的包含映射,其微分在任一光滑坐标下恰好是把切向量的分量以恒等方式映到较大环境的切向量中。下面的命题表明,把 $T_pS$ 视为 $T_pM$ 的子空间是自然的,即在 $S$ 中任意光滑曲线的切向量可被识别为它在 $M$ 中的切向量。

\begin{proposition}\label{proposition:ch05-tangent-vector-identification}
设 $S\subseteq M$ 是浸入子流形, $p\in S$。在上述等同下, $T_pS$ 恰好由 $M$ 中切于 $S$ 的那些向量组成,即由形如 $\gamma'(0)$ 的向量组成,其中 $\gamma:(-\varepsilon,\varepsilon)\to S$ 是 $S$ 中的光滑曲线,看作 $M$ 中的曲线。
\end{proposition}

\begin{proof}
若 $\gamma:(-\varepsilon,\varepsilon)\to S$ 是 $S$ 中的光滑曲线,则 $\gamma$ 看作 $M$ 中的曲线就是 $\iota\circ\gamma$,故其在 $0$ 处的切向量为 $\mathrm{d}\iota_p(\gamma'(0))$,它位于 $\mathrm{d}\iota_p(T_pS)$ 中。反之,任一向量 $v\in T_pS$ 可写为 $v=\gamma'(0)$,其中 $\gamma$ 是 $S$ 中的光滑曲线,故 $\mathrm{d}\iota_p(v)$ 是曲线 $\iota\circ\gamma$ 在 $M$ 中的切向量。
\end{proof}

在切片坐标下, $S$ 对应于 $x^{k+1}=\cdots=x^n=0$,而包含映射在坐标下由 $(x^1,\dots,x^k)\mapsto(x^1,\dots,x^k,0,\dots,0)$ 给出。因此,把 $T_pS$ 等同于 $T_pM$ 的子空间,就恰好是把它等同于由 $\partial/\partial x^1|_p,\dots,\partial/\partial x^k|_p$ 张成的子空间。下面的定理表明,在嵌入子流形的情形,这与用嵌入子流形 $S$ 的任一局部定义函数来刻画切空间是一致的。

\begin{theorem}[The Tangent Space to a Level Set]\label{theorem:ch05-tangent-level-set}
设 $S\subseteq M$ 是嵌入子流形, $\Phi:U\to N$ 是 $S$ 在开子集 $U\subseteq M$ 上的局部定义函数。对任一 $p\in S\cap U$, $T_pS\subseteq T_pM$ 恰为映射 $\mathrm{d}\Phi_p:T_pM\to T_{\Phi(p)}N$ 的核。
\end{theorem}

\begin{proof}
由切片坐标的存在性,可选取以 $p$ 为中心的 $M$ 的切片坐标 $(x^1,\dots,x^n)$,使 $S$ 局部由 $x^{k+1}=\cdots=x^n=0$ 给出,其中 $k=\dim S$。我们也可选取 $\Phi$ 的坐标表示具有形式\eqref{align:ch04-rank-canonical},因为它是局部定义函数,即满射。于是 $\mathrm{d}\Phi_p$ 的核恰好由 $\partial/\partial x^1|_p,\dots,\partial/\partial x^k|_p$ 张成,而这正是 $T_pS$。
\end{proof}

\begin{example}[The Tangent Space to a Sphere]\label{example:ch05-tangent-sphere}
利用 $S^n\subseteq\mathbb{R}^{n+1}$ 是 $f(x)=|x|^2$ 的正则水平集这一事实,上一定理表明 $S^n$ 在点 $a\in S^n$ 的切空间 $T_aS^n$ 是 $\mathrm{d}f_a:T_a\mathbb{R}^{n+1}\to\mathbb{R}$ 的核。若我们用 $T_a\mathbb{R}^{n+1}$ 等同于 $\mathbb{R}^{n+1}$,则 $\mathrm{d}f_a(v)=2\sum_i a^iv^i=2\langle a,v\rangle$,其中 $\langle a,v\rangle$ 是通常的 Euclid 内积。于是 $T_aS^n=\{v:\langle a,v\rangle=0\}$,即与 $a$ 正交的向量组成的 $n$ 维子空间。
\end{example}

\begin{figure}[ht]
\centering
\includegraphics[width=0.8\textwidth]{fig-05-12.png}
\caption{嵌入子流形的切空间}\label{figure:ch05-tangent-embedded}
\end{figure}

上面关于切空间到 $M$ 的切空间的识别,有一个简单而重要的推论。设 $M$ 是带边光滑流形。\cref{theorem:ch05-boundary-submanifold} 表明 $\partial M$ 是 $M$ 的真嵌入子流形,故我们可把 $T_p(\partial M)$ 等同于 $T_pM$ 的一个 $n-1$ 维子空间。回忆在边界点处, $M$ 的切空间与全空间 $T_pM$ 等同。边界的切空间是其中的一个超平面。

对 $M$ 的边界点 $p$,称 $M$ 中的非零切向量 $v\in T_pM$ 指向外,如果它可表为 $\gamma'(0)$,其中 $\gamma:(-\varepsilon,\varepsilon)\to M$ 是 $M$ 中的光滑曲线,且对 $t>0$ 充分小,有 $\gamma(t)\notin M$(\cref{figure:ch05-outward-vector})。

\begin{figure}[ht]
\centering
\includegraphics[width=0.8\textwidth]{fig-05-9.png}
\caption{限制定义域}\label{figure:ch05-outward-vector}
\end{figure}

\begin{proposition}[The Tangent Space to a Boundary]\label{proposition:ch05-tangent-boundary}
设 $M$ 是带边光滑 $n$ 流形, $p\in\partial M$,并设 $(x^1,\dots,x^n)$ 是 $M$ 上以 $p$ 为中心的任一光滑边界坐标卡。
\begin{enumerate}[(a)]
\item $T_p(\partial M)$ 恰为 $\partial/\partial x^1|_p,\dots,\partial/\partial x^{n-1}|_p$ 张成的 $T_pM$ 的 $(n-1)$ 维子空间。
\item 向量 $v=v^i\partial/\partial x^i|_p\in T_pM$ 指向外,当且仅当 $v^n<0$。
\end{enumerate}
\end{proposition}

\begin{proof}
见习题 5-10。
\end{proof}

\begin{figure}[ht]
\centering
\includegraphics[width=0.8\textwidth]{fig-05-13.png}
\caption{内向向量}\label{figure:ch05-inward-vector}
\end{figure}

当 $S\subseteq M$ 是嵌入子流形时,其余维数 $k=\dim M-\dim S$ 有一个重要解释:在任一点 $p\in S$, $T_pS$ 是 $T_pM$ 的 $k$ 余维子空间。若 $M$ 是 Riemannian 流形(第十二章),则有一个自然的正交补,称为 $S$ 在 $p$ 处的法空间,并记为 $N_pS=\bigl(T_pS\bigr)^\perp$,它由所有与 $T_pS$ 中每个向量正交的向量组成。然而,在没有度量的情况下,没有自然的方式把法方向从 $T_pM$ 中分离出来。

\end{document}